% Schede riscontrate sul codice corrente; non costituiscono esiti di test.
\newcommand{\scheda}[7]{%
  \subsection{\texorpdfstring{\raggedright #1 --- #2}{#1 --- #2}}\label{uc:#1}
  \begin{longtable}{|p{.20\linewidth}|p{.72\linewidth}|}
  \caption{Scheda del caso d'uso #1}\\\hline
  \textbf{Campo} & \textbf{Contenuto}\\\hline\endfirsthead
  \textbf{Campo} & \textbf{Contenuto (continua)}\\\hline\endhead
  Identificativo & #1\\\hline
  Descrizione & #3\\\hline
  Livello & #4\\\hline
  Attori & #5\\\hline
  Precondizioni & #6\\\hline
  Postcondizioni & #7\\\hline
  \end{longtable}
}
\newcommand{\principale}{\paragraph{Flusso principale}}
\newcommand{\alternative}{\paragraph{Alternative}}

Ogni titolo riporta identificativo stabile e denominazione. Le alternative
\emph{n.a}, \emph{n.b} si diramano dal passo \emph{n} della stessa scheda;
si indica se il caso riprende o termina. O indica un obiettivo applicativo,
F una sottofunzione applicativa e S una funzione di supporto. Sono livelli
descrittivi, non livelli di autorizzazione. Gli attori sono ruoli dell'interfaccia;
non si assume un controllo dei permessi indipendente in ogni servizio.
Le postcondizioni descrivono il successo; in caso di errore transazionale
le scritture della singola operazione vengono annullate. Un errore successivo
nell'aggiornamento della vista non annulla un'operazione gi\`a confermata.


Le prove richiamate sono sorgenti e risultati storici, non nuove esecuzioni.
\clearpage
\scheda{UC-01}{Prenotare un test drive}
{Il cliente riserva un giro di prova di un'ora su un veicolo della concessionaria.}
{O --- obiettivo applicativo.}
{Cliente; indirettamente altri clienti con prenotazioni concorrenti.}
{Cliente autenticato; veicolo visibile nel catalogo vendita o noleggio.}
{Prenotazione associata a cliente, veicolo e data/ora; nessun addebito.}
\principale
\begin{enumerate}
\item Il cliente seleziona un veicolo in \emph{Auto in Vendita} o \emph{Auto a Noleggio}.
\item Apre \emph{Prenota Test Drive}, sceglie data e ora iniziale e conferma.
      La UI propone ore intere dalle 09:00 alle 16:00.
\item Il servizio rilegge il veicolo, ammette solo \code{SALE}/\code{RENT},
      controlla che l'istante non sia passato, che il cliente abbia meno di tre
      test drive attivi e che slot e giorno non siano occupati da test drive o noleggi incompatibili.
\item Il sistema salva la prenotazione nella transazione e aggiorna \emph{Le Mie Operazioni}.
\item Il cliente legge la conferma con data e ora del giro di prova.
\end{enumerate}
\alternative
\begin{description}
\item[2.a] Annullamento del dialogo: nessuna prenotazione; termina.
\item[3.a] Data passata, stato cambiato, limite raggiunto o indisponibilit\`a:
  errore, nessuna nuova prenotazione; il cliente pu\`o ripartire da 1 o 2.
\item[4.a] Errore di persistenza: rollback della prenotazione; termina senza successo.
\item[5.a] Successivamente il cliente seleziona un proprio test drive in
  \emph{Le Mie Operazioni} e lo annulla: viene rimosso e la tabella aggiornata.
  Un identificativo non appartenente al cliente viene rifiutato.
\item[5.b] Un successivo noleggio o completamento di vendita sul veicolo
  cancella tutti i suoi test drive futuri, anche fuori dal periodo noleggiato.
  Lo storico rimane; la cancellazione partecipa alla transazione commerciale.
\end{description}
\begin{flushleft}\textbf{Riscontro:} \code{TestDriveService}, \code{TestDriveCancellationObserver},
\code{ClientDashboardView}. La durata \`e una convenzione del giro di prova;
la UI seleziona ore intere, senza attribuire al servizio un vincolo generale
sui minuti che il codice non applica per tutte le ore.\end{flushleft}


\begin{flushleft}\noindent\textbf{Prove e portata:} \code{BookingRulesTest} (fascia e date), \code{BookingConcurrencyIntegrationTest} (limite, slot concorrenti e rollback), \code{TestDriveCancellationIntegrationTest} (cancellazioni); evidenze storiche \code{r06-bookings}/\code{r11-integrated}.\end{flushleft}
\clearpage
\scheda{UC-02}{Noleggiare un veicolo}
{La concessionaria prenota un proprio veicolo al cliente e incassa subito il costo del periodo.}
{O --- obiettivo applicativo.}
{Cliente; clienti titolari di test drive futuri sul medesimo veicolo.}
{Cliente autenticato con conto demo; veicolo presente nel catalogo noleggio.}
{Noleggio salvato e costo intero addebitato; test drive futuri del veicolo cancellati;
il veicolo conserva la destinazione \code{RENT}.}
\principale
\begin{enumerate}
\item Il cliente seleziona un veicolo in \emph{Auto a Noleggio} e apre \emph{Prenota Noleggio}.
\item Sceglie inizio e fine; il dialogo mostra il costo del periodo con estremi inclusi.
\item Conferma. Il servizio rilegge stato, disponibilit\`a e prezzo: inizio non passato,
      fine non precedente all'inizio, stato \code{RENT}, nessuna sovrapposizione.
\item Il sistema verifica il conto e addebita il 100\% del costo ricalcolato.
\item Salva il noleggio e cancella tutti i test drive futuri del veicolo,
      anche oltre il periodo prenotato, nella stessa transazione.
\item Mostra la conferma e aggiorna saldo e operazioni del cliente.
\end{enumerate}
\alternative
\begin{description}
\item[2.a] Date incomplete o invertite: il preventivo non \`e valido;
  il cliente corregge il periodo al passo 2 oppure annulla senza effetti.
\item[3.a] Stato cambiato o periodo gi\`a occupato, anche per una prenotazione concorrente:
  rifiuto senza addebito; ripartire da 1 o 2.
\item[4.a] Conto assente o saldo insufficiente: nessun noleggio e nessun addebito; termina.
\item[5.a] Errore nel salvataggio o nelle cancellazioni collegate: rollback
  anche dell'addebito; nessun successo parziale.
\item[6.a] In seguito UC-08 destina il veicolo alla vendita: se non vi sono
  noleggi in corso, quelli con inizio futuro sono rimossi e rimborsati al 100\%.
  Il cliente vede l'effetto ricaricando le proprie operazioni.
\end{description}
\begin{flushleft}\textbf{Riscontro:} \code{RentalService}, \code{PriceCalculatorService},
\code{TestDriveCancellationObserver}. Il preventivo non \`e un accordo fissato:
il costo viene ricalcolato alla prenotazione. La UI non offre una cancellazione
autonoma del noleggio al cliente.\end{flushleft}


\begin{flushleft}\noindent\textbf{Prove e portata:} \code{BookingConcurrencyIntegrationTest} (conflitti e addebito), \code{TestDriveCancellationIntegrationTest} (rollback Observer), \code{VehicleDestinationIntegrationTest} (rimborsi); evidenze \code{r06-bookings}/\code{r06-destination}/\code{r11-integrated}.\end{flushleft}
\clearpage
\scheda{UC-03}{Richiedere e completare una vendita al cliente}
{La concessionaria vende un veicolo: acconto, autorizzazione del venditore e saldo del totale concordato.}
{O --- obiettivo applicativo articolato in pi\`u sessioni.}
{Cliente principale; venditore per UC-04.}
{Cliente autenticato con conto demo; veicolo proposto in vendita.}
{Nel successo completo: totale registrato come vendita al cliente,
richiesta rimossa, veicolo \code{SOLD}; test drive futuri cancellati.
Dopo il solo acconto: richiesta \code{SALE} aperta, totale fissato e veicolo \code{BLOCKED}.}
\principale
\begin{enumerate}
\item Il cliente seleziona il veicolo in \emph{Auto in Vendita} e avvia l'acquisto.
\item Legge prezzo indicativo e acconto del 20\% nel dialogo e conferma.
\item Il servizio riserva il veicolo solo se ancora \code{SALE}, ricalcola il prezzo,
  addebita l'acconto e salva la richiesta con totale concordato, senza venditore n\'e scadenza.
\item Il venditore approva la richiesta attraverso UC-04, fissando il termine di un mese.
\item Il cliente seleziona la richiesta approvata in \emph{Le Mie Operazioni},
  apre il pagamento del saldo e conferma il residuo del totale concordato.
\item Il servizio verifica autorizzazione e scadenza, addebita il residuo,
  registra l'intero totale nel registro, elimina la richiesta e porta il veicolo a \code{SOLD}.
  Cancella i test drive futuri nella stessa transazione del saldo.
\end{enumerate}
\alternative
\begin{description}
\item[2.a] Annullamento: il veicolo non viene riservato; termina.
\item[3.a] Veicolo non pi\`u disponibile, conto assente o fondi insufficienti:
  rollback di acconto, richiesta e stato; termina.
\item[4.a] Rifiuto del venditore: rimborso integrale dell'acconto,
  richiesta eliminata e veicolo di nuovo \code{SALE}; termina.
\item[5.a] Richiesta non autorizzata o gi\`a scaduta: il saldo viene rifiutato.
  Annullare il dialogo lascia invece aperta la richiesta.
\item[6.a] Saldo insufficiente o errore transazionale: nessuna vendita completata;
  restano acconto e richiesta precedenti; il cliente pu\`o riprovare entro il termine.
\item[6.b] Mancato saldo entro il termine: all'accesso successivo il controllo
  delle scadenze trattiene il 20\% dell'acconto come penale, restituisce il resto,
  rimuove la richiesta e libera il veicolo per la vendita. Non \`e un processo periodico in background.
\end{description}
\begin{flushleft}\textbf{Riscontro:} \code{PurchaseService}, \code{Request}, ADR 0001.
Acconto e saldo hanno transazioni distinte. Basi e sconti successivi non
modificano il totale gi\`a concordato; la penale \`e il 20\% dell'acconto, non del totale.\end{flushleft}


\begin{flushleft}\noindent\textbf{Prove e portata:} \code{PurchaseAgreementIntegrationTest} (accordo immutato, autorizzazione, scadenza, fondi e rollback), \code{TestDriveCancellationIntegrationTest} (Observer); evidenze \code{r04-agreement}/\code{r11-integrated}.\end{flushleft}
\clearpage
\scheda{UC-05}{Proporre un veicolo alla concessionaria}
{La concessionaria valuta e acquista il veicolo ceduto dal cliente; l'etichetta UI \emph{Vendi la tua Auto} usa il punto di vista del cliente.}
{O --- obiettivo applicativo in due fasi, proposta ed esito.}
{Cliente proponente; venditore che decide.}
{Cliente autenticato; modello gi\`a censito; dati del veicolo disponibili.
Se il VIN \`e noto, stato \code{SOLD}, ultima vendita al proponente e nessuna richiesta aperta.}
{Accettazione: stima accreditata al cliente, acquisto della concessionaria registrato,
richiesta rimossa e veicolo \code{SALE}. Dopo la proposta: veicolo \code{WAITING} e richiesta \code{PURCHASE}.}
\principale
\begin{enumerate}
\item Il cliente apre \emph{Vendi la tua Auto} e indica marchio, modello,
  VIN, targa, chilometri e anno.
\item Pu\`o calcolare il preventivo, basato sul modello e sui dati inseriti.
\item Invia la proposta. Il servizio valida i dati e l'eventuale titolarit\`a
  del VIN noto; crea il veicolo oppure aggiorna quello gi\`a venduto al cliente.
\item Salva il veicolo \code{WAITING} e la richiesta con stima ricalcolata,
  positiva o nulla, senza effettuare pagamenti.
\item Il venditore seleziona la proposta \emph{Acquisto dal cliente} in
  \emph{Richieste in Sospeso} e aziona \emph{Approva Richiesta}.
\item Il sistema accredita la stima memorizzata, registra l'acquisto dal cliente
  con venditore, porta il veicolo a \code{SALE} ed elimina la richiesta,
  nella stessa transazione. La UI mostra l'esito.
\end{enumerate}
\alternative
\begin{description}
\item[2.a] Preventivo non richiesto: il cliente pu\`o passare direttamente a 3;
  il servizio ricalcola comunque la stima all'invio.
\item[3.a] Dati invalidi, anno precedente al modello, VIN occupato o non
  appartenente al proponente, richiesta gi\`a aperta: errore e nessuna nuova proposta;
  correggere al passo 1.
\item[5.a] Rifiuto: nessun accredito; richiesta eliminata.
  Il veicolo senza operazioni associate viene eliminato; con storico torna \code{SOLD}.
\item[6.a] Stima nulla, conto mancante, richiesta non pi\`u valida o errore di persistenza:
  rollback dell'intero esito, senza pagamento o acquisizione parziali.
\end{description}
\begin{flushleft}\textbf{Riscontro:} \code{SaleToDealershipService},
\code{ClientDashboardView}, \code{SalesmanDashboardView}. Il registro conserva
un importo non negativo con tipo acquisto; il suo significato \`e un esborso
per la concessionaria. Non esiste un successivo pagamento a saldo per la cessione.
La proposta ammette una stima nulla, ma il conto rifiuta l'accredito zero:
l'accettazione fallisce con rollback (sezione~\ref{sec:sequenze}).\end{flushleft}


\begin{flushleft}\noindent\textbf{Prove e portata:} \code{SaleProposalIntegrationTest} (titolarità e VIN), \code{SaleOutcomeIntegrationTest} (accredito, storico, concorrenza e rollback); evidenze \code{r05-sale-outcome}/\code{r11-integrated}.\end{flushleft}
\clearpage
\scheda{UC-04}{Autorizzare una richiesta di vendita al cliente}
{Il venditore decide l'esito della richiesta con cui la concessionaria vender\`a un veicolo al cliente.}
{F --- sottofunzione applicativa di UC-03, svolta dal venditore.}
{Venditore principale; cliente destinatario dell'esito.}
{Venditore autenticato; richiesta \code{SALE} in sospeso, con acconto gi\`a versato e veicolo riservato.}
{Approvazione: richiesta attribuita al venditore e scadenza impostata a un mese da oggi;
il saldo resta da pagare e il veicolo rimane \code{BLOCKED}.}
\principale
\begin{enumerate}
\item Il venditore apre \emph{Richieste in Sospeso}.
\item Seleziona una riga \emph{Vendita al cliente}, verificando totale,
  acconto e residuo mostrati nella tabella.
\item Aziona \emph{Approva Richiesta}; non compare una conferma preventiva.
\item Il servizio rilegge la richiesta, verifica che sia \code{SALE},
  assegna il venditore e la scadenza a un mese dalla data corrente.
\item La UI mostra l'esito e ricarica richieste e inventario;
  il cliente potr\`a saldare in UC-03.
\end{enumerate}
\alternative
\begin{description}
\item[2.a] Nessuna riga selezionata: messaggio di selezione mancante; ritorno a 2.
\item[2.b] Riga \emph{Acquisto dal cliente}: non \`e questa autorizzazione;
  la stessa azione conclude la cessione secondo UC-05.
\item[3.a] \emph{Rifiuta Richiesta}: il servizio restituisce l'intero acconto,
  elimina la richiesta e ripristina \code{SALE} nella stessa transazione.
  La vendita non prosegue; nessuna penale.
\item[4.a] Richiesta rimossa o tipo non valido: errore, nessuna nuova autorizzazione;
  il venditore ricarica l'elenco.
\item[4.b] Errore durante l'aggiornamento oppure durante il rimborso di 3.a:
  rollback della decisione; lo stato precedente rimane valido.
\end{description}
\begin{flushleft}\textbf{Riscontro:} \code{SalesmanDashboardView},
\code{PurchaseService.approvePurchase}, \code{PurchaseService.rejectPurchase},
\code{OperationQueryService}. Il caso non registra un nuovo venditore e non
rappresenta l'abilitazione di un account staff. La selezione delle richieste
in sospeso \`e affidata alla query della UI; non si attribuisce al comando
un controllo aggiuntivo contro ogni possibile riapprovazione da un chiamante diverso.\end{flushleft}


\begin{flushleft}\noindent\textbf{Prove e portata:} \code{PurchaseAgreementIntegrationTest} (approvazione e termine), \code{PurchaseConcurrencyIntegrationTest} (rifiuto, rimborso ed esiti concorrenti); l\'assenza di selezione e la riapprovazione non sono certificate da quel test. Evidenze \code{r04-agreement}/\code{r11-integrated}.\end{flushleft}
\clearpage
\scheda{UC-06}{Gestire lo sconto su un veicolo}
{Il manager applica o rimuove una promozione di modello o di singolo veicolo.}
{O --- obiettivo applicativo di gestione commerciale.}
{Manager; clienti interessati ai nuovi prezzi.}
{Manager autenticato; modello o veicolo destinatario gi\`a presente nel catalogo.}
{Sconto salvato o rimosso e tabella aggiornata. I nuovi calcoli usano il massimo
fra sconto modello e veicolo; i totali concordati nelle richieste aperte restano invariati.}
\principale
\begin{enumerate}
\item Il manager apre la scheda \emph{Sconti}.
\item Nel riquadro modello seleziona un modello; nel riquadro veicolo
  inserisce invece l'identificativo del veicolo.
\item Inserisce una percentuale intera tra 0 e 50 e aziona \emph{Salva} nel riquadro scelto.
\item Il servizio valida la percentuale e persiste lo sconto;
  uno sconto gi\`a presente sullo stesso elemento viene aggiornato.
\item La UI mostra la conferma e ricarica le due tabelle degli sconti.
\end{enumerate}
\alternative
\begin{description}
\item[2.a] Modello non selezionato o ID non numerico: errore prima del salvataggio;
  ritorno a 2. Un ID veicolo inesistente viene rifiutato dalla persistenza.
\item[3.a] Percentuale non intera o esterna a 0--50: errore, nessuna modifica; ritorno a 3.
\item[3.b] Il manager seleziona uno sconto esistente e aziona \emph{Rimuovi}:
  il servizio elimina lo sconto e la tabella viene ricaricata.
  Senza selezione l'azione non modifica dati; non compare una conferma preventiva.
\item[4.a] Errore transazionale: nessun nuovo sconto salvato; termina senza successo.
\item[5.a] Un'altra tabella gi\`a aperta mostra ancora il vecchio prezzo:
  ricaricarla con \emph{Aggiorna} oppure riaprire la vista.
  La UI non propaga automaticamente ogni modifica a tutte le dashboard gi\`a aperte.
\end{description}
\begin{flushleft}\textbf{Riscontro:} \code{CatalogService}, \code{Discount},
\code{PriceCalculatorService}, \code{ManagerDashboardView}. Gli sconti non
si sommano. La base manuale del veicolo sostituisce la base del modello nella
stessa formula; non impone un prezzo finale sottratto agli sconti.\end{flushleft}


\begin{flushleft}\noindent\textbf{Prove e portata:} \code{PriceCalculatorServiceTest} (massimo e base manuale), \code{DomainInvariantsTest} (0--50), \code{PriceSelectionIntegrationTest} (persistenza della base); evidenze \code{r03-prices}/\code{r11-integrated}.\end{flushleft}
